\documentclass[11pt]{article}
\usepackage{amsmath,amssymb,amsthm}
\usepackage[margin=1in]{geometry}

\newtheorem{theorem}{Theorem}
\newtheorem{corollary}[theorem]{Corollary}
\newtheorem{proposition}[theorem]{Proposition}

\title{V82: Hall-Neighborhood Minima and Transversal Girth}
\author{NC0\_k-Avoid Laboratory}
\date{}

\begin{document}
\maketitle

Let \(G=(M,X;E)\) be a bipartite graph and let \(T(G)\) be the
transversal matroid on \(M\). For \(S\subseteq M\), write \(N(S)\) for
its neighborhood and define
\[
  h^*(G)=\min\{|N(S)|:|N(S)|<|S|\}.
\]
We assume that the minimum is defined, equivalently that \(T(G)\) has a
dependent set.

\begin{theorem}[Hall--girth equivalence]
\[
  h^*(G)=\operatorname{girth}(T(G))-1.
\]
\end{theorem}

\begin{proof}
Choose an inclusion-minimal set \(S\) attaining \(h^*(G)\). If a
proper subset \(R\subsetneq S\) were Hall deficient, then
\(N(R)\subseteq N(S)\), hence \(|N(R)|\le h^*(G)\). Minimality of the
objective forces equality, contradicting the inclusion-minimal choice
of \(S\). Therefore \(S\) is a circuit of \(T(G)\).

For each \(e\in S\), the set \(S-\{e\}\) has a matching into \(N(S)\);
thus \(|N(S)|\ge |S|-1\). Since \(S\) is dependent,
\(|N(S)|<|S|\), and consequently \(|N(S)|=|S|-1\). This gives
\(h^*(G)\ge \operatorname{girth}(T(G))-1\).

Conversely, a shortest circuit \(C\) satisfies the same equality
\(|N(C)|=|C|-1\), yielding the reverse inequality.
\end{proof}

\begin{corollary}[Deficiency-one separation]
Every inclusion-minimal minimizer of \(h^*(G)\) has Hall deficiency
exactly one.
\end{corollary}

Thus minimum neighborhood and maximum deficiency are distinct
optimization objectives. The former yields a circuit; the latter can
improve random-sampling success but need not minimize the number of
variables.

\begin{proposition}[Degree-two bicircular boundary]
If every element of \(M\) has at most two neighbors in \(X\), then
\(T(G)\) is the bicircular matroid of the multigraph obtained by turning
each two-element support into an edge and each singleton support into a
loop. With this graph presentation supplied, its girth is computable
in polynomial time.
\end{proposition}

\begin{proof}
A set of edges admits an assignment to distinct incident vertices
exactly when each connected component is a pseudoforest. Hence the
circuits are the minimal connected subgraphs with two cycles: theta
graphs, tight handcuffs, and loose handcuffs. For each of these three
fixed topologies, enumerate its constant number of branch or attachment
vertices and solve the resulting fixed-terminal disjoint-path/minimum-
cost-flow instance. Taking the minimum over all enumerations gives a
polynomial algorithm.
\end{proof}

The published general hardness theorem of Colbourn and Elmallah
(\emph{Discrete Mathematics} 114 (1993), Theorem 2.1) implies, through
the first theorem, that exact \(h^*\) is NP-hard for unrestricted
transversal presentations. The exact complexity when every element
has at most three neighbors is not resolved here.

\paragraph{Nonclaims.}
This module proves no degree-three hardness or algorithm, no circuit
lower bound, and no result resolving P versus NP.

\end{document}
